% Part: incompleteness
% Chapter: arithmetization-syntax
% Section: coding-symbols

\documentclass[../../../include/open-logic-section]{subfiles}

\begin{document}

% SEGMENT OLP-0282-S01

\olfileid{inc}{art}{cod}
\olsection{குறிகளைக் குறியீடாக்குதல்}

முதல்தர தருக்கத்தின் அடிப்படை மொழி~$\Lang L$ பின்வரும் குறிகளைப்
பயன்படுத்துகிறது:
\[
\lfalse \quad \lnot \quad \lor \quad \land \quad \lif \quad \lforall
\quad \lexists \quad \eq \quad ( \quad ) \quad ,
\]
இவற்றுடன் மாறிகள், !!{constant}s ஆகியவற்றின் !!{enumerable} கணங்களும்,
தன்னிச்சையான இடவெண் கொண்ட !!{function}s, !!{predicate}s ஆகியவற்றின்
!!{enumerable} கணங்களும் பயன்படுகின்றன. ஒவ்வொரு குறிக்கும் தனித்துவமான
ஓர் எண்ணைக் \emph{குறியீடாக} ஒதுக்கி, வெவ்வேறு குறிகள் இரண்டுக்கு ஒரே
எண் ஒதுக்கப்படாதவாறு இக்குறிகள் ஒவ்வொன்றுக்கும் குறியீடுகளை வழங்கலாம்.
எல்லாக் குறிகளின் கணமும் !!{enumerable} என்பதாலும், அதற்கும் இயல் எண்களின்
கணத்திற்கும் இடையே !!a{bijection} இருப்பதாலும் இது சாத்தியம் என்பதை
அறிவோம். ஆனால் ஒரு குறியீட்டிலிருந்து குறியையும் (அதைப் பற்றிய சில தகவலையும்,
எடுத்துக்காட்டாக !!a{function} இன் இடவெண்ணையும்) கணிக்கத்தக்க வழியில்
மீட்க முடிவதை உறுதிசெய்ய விரும்புகிறோம். இதைச் செய்வதற்குப் பல வழிகள்
உள்ளன. முதன்மை மீள்வரையறைச் சார்புகளைப் பயன்படுத்தும் அத்தகைய ஒரு வழி
இதோ. ($\tuple{n_0, \dots, n_k}$ என்பது $n_0$, \dots, $n_k$ என்ற
எண்களின் தொடரைக் குறியீடாக்கும் எண் என்பதை நினைவுகூர்க.)

\begin{defn}
$s$ ஆனது~$\Lang L$ இன் ஒரு குறி எனில், அதன் \emph{குறிக்
குறியீடு}~$\scode s$ பின்வருமாறு வரையறுக்கப்படுகிறது:
\begin{enumerate}
\item $s$ தருக்கக் குறிகளுள் ஒன்று எனில், பின்வரும் அட்டவணை~$\scode s$
  ஐத் தருகிறது:
\[
\begin{array}{cccccccccc}
  \lfalse & \lnot & \lor & \land & \lif & \lforall \\
\tuple{0, 0} & \tuple{0, 1} & \tuple{0, 2} & \tuple{0, 3} &
\tuple{0, 4} & \tuple{0, 5} \\
\lexists & \eq & ( & ) & ,\\
\tuple{0, 6} & \tuple{0, 7} &
\tuple{0, 8} & \tuple{0, 9} & \tuple{0, 10}
\end{array}
\]
\item $s$ ஆனது $i$-ஆம் மாறி~$\Obj v_i$ எனில், $\scode s = \tuple{1, i}$.
\item $s$ ஆனது $i$-ஆம் !!{constant}~$\Obj c_i$ எனில்,
  $\scode s = \tuple{2, i}$.
\item $s$ ஆனது $i$-ஆம் $n$-இட !!{function}~$\Obj f_i^n$ எனில்,
  $\scode s = \tuple{3, n, i}$.
\item $s$ ஆனது $i$-ஆம் $n$-இட !!{predicate}~$\Obj P_i^n$ எனில்,
  $\scode s = \tuple{4, n, i}$.
\end{enumerate}
\end{defn}

\begin{prop}
பின்வரும் தொடர்புகள் முதன்மை மீள்வரையறையானவை:
\begin{enumerate}
\item ஏதாவது~$i$ க்கு $x$ ஆனது~$\Obj f^n_i$ இன் குறியீடாக, அதாவது $x$
  ஓர் $n$-இட !!{function} இன் குறியீடாக இருந்தால், அப்போதும் அப்போது மட்டுமே
  $\fn{Fn}(x, n)$.
\item ஏதாவது~$i$ க்கு $x$ ஆனது~$\Obj P^n_i$ இன் குறியீடாகவோ, அல்லது
  $x$ ஆனது~$\eq$ இன் குறியீடாகவும் $n = 2$ ஆகவோ இருந்தால், அதாவது $x$
  ஓர் $n$-இட !!{predicate} இன் குறியீடாக இருந்தால், அப்போதும் அப்போது
  மட்டுமே $\fn{Pred}(x, n)$.
\end{enumerate}
\end{prop}

\begin{defn}
$s_0, \dots, s_{n-1}$ என்பது குறிகளின் தொடர் எனில், அதன் \emph{கோடல்
  எண்} $\tuple{\scode{s_0}, \dots, \scode{s_{n-1}}}$ ஆகும்.
\end{defn}

\begin{explain}
\emph{குறியீடுகளும்} \emph{கோடல் எண்களும்} வெவ்வேறு பொருள்கள் என்பதைக்
கவனிக்கவும். எடுத்துக்காட்டாக, மாறி~$\Obj v_5$ க்கு
$\scode{\Obj v_5} = \tuple{1, 5} = 2^2\cdot 3^6$ என்ற குறியீடு உள்ளது.
ஆனால் ஒரு சொல்லாகக் கருதப்படும் மாறி~$\Obj v_5$, (நீளம்~$1$ கொண்ட)
குறிகளின் தொடராகவும் உள்ளது. \emph{சொல்}~$\Obj v_5$ இன் \emph{கோடல்
எண்}~$\Gn{\Obj v_5}$ என்பது $\tuple{\scode{\Obj v_5}} = 2^{\scode{\Obj
    v_5} + 1} = 2^{2^2\cdot 3^6 + 1}$.
\end{explain}

\begin{ex}
$k_0$, \dots, $k_{n-1}$ என்பது எண்களின் தொடர் எனில், பகா எண் அடுக்குக்
குறியீடாக்கத்தில் அத்தொடரின் குறியீடு~$\tuple{k_0, \dots, k_{n-1}}$
பின்வருவது என்பதை நினைவுகூர்க:
\[
2^{k_0+1}\cdot3^{k_1+1}\cdot \dots \cdot p_{n-1}^{k_{n-1}+1},
\]
இங்கு $p_i$ என்பது $i$-ஆம் பகா எண் ($p_0 = 2$ இல் தொடங்குகிறது).
ஆகவே, எடுத்துக்காட்டாக, வாய்பாடு $\eq[\Obj v_0][\Obj 0]$, அல்லது இன்னும்
வெளிப்படையாக ${\eq}(\Obj v_0,\Obj c_0)$, பின்வரும் கோடல் எண்ணைக் கொண்டது:
\[
\tuple{\scode{\eq},\scode{(},\scode{\Obj v_0},\scode{,}, \scode{\Obj
    c_0},\scode{)}}.
\]
இங்கு $\scode{\eq}$ என்பது $\tuple{0,7} = 2^{0+1}\cdot
3^{7+1}$; $\scode{\Obj v_0}$ என்பது $\tuple{1,0} = 2^{1+1}\cdot3^{0+1}$;
இவ்வாறே மற்றவையும் அமையும். ஆகவே $\Gn{=(\Obj v_0,\Obj c_0)}$ என்பது
\begin{multline*}
2^{\scode{=} + 1}\cdot 3^{\scode{(}+1}\cdot 5^{\scode{\Obj v_0}+1}
\cdot 7^{\scode{,} + 1} \cdot 11^{\scode{\Obj c_0}+1} \cdot
13^{\scode{)}+1} = \\
2^{2^1\cdot 3^8 + 1}\cdot 3^{2^1\cdot 3^9+1}\cdot 5^{2^2\cdot 3^1+1}
\cdot 7^{2^1\cdot 3^{11} + 1} \cdot 11^{2^3\cdot3^1+1} \cdot
13^{2^1\cdot3^{10}+1} = \\
2^{13\,123}\cdot 3^{39\,367}\cdot 5^{13}\cdot 7^{354\,295}\cdot11^{25}\cdot13^{118\,099}.
\end{multline*}
\end{ex}

\end{document}
